% Section 8 --- What is missing
% normalisation needed (no smart-quotes, no underscores). Roadmap section;
% introduces 3 new citations: Witman2018JCTC, Becker2018JPCC, McCready2024JCTC.
\section{What is missing}\label{sec:missing}

The v0.4 result is deliberately incomplete. It proves that two locked recipes reproduce experiment within the strict band, and that three further recipes are physically credible within the wider Tier-B band. It does not prove that the atlas is finished. The value of the audit is that the remaining gaps are now named scientific tasks rather than hidden sources of error.

The first missing piece is a method for flexible or gate-opening adsorption. Na-Rho is the clearest example. A rigid closed-state Widom calculation cannot become a scalar validation of an experiment that involves a trapdoor opening process. Closing this gap requires a method that samples the relevant framework state or transition coordinate, not a different label on the same closed-state calculation. Possible routes include flat-histogram or transition-matrix approaches, or a mechanism-aware open-state adsorption calculation, but those would be new calculations rather than a reinterpretation of the v0.4 branch \cite{Witman2018JCTC}.

The second missing piece is a completed cationic-zeolite input set. The 5c family is already useful as a reference audit, but it is not yet an executed validation family. Na-ZK-5, Zeolite 5A, Zeolite 13X, and Zeolite 4A require cation-containing structures, cation positions, cation charges, and locked cation force-field parameters before the corresponding Widom recipes are defined. A topology-only or pure-silica framework would test the wrong material. The next step is therefore not simply to run more insertions, but to lock the correct structures and cation models.

The third missing piece is a better treatment of strong open-metal sites. The Mg-MOF-74 failures show that ordinary non-polarisable or simplified pairwise recipes are not enough for exposed metal centres (HKUST-1 under genuine UFF, by contrast, reaches the Tier-B band once its open-Cu site is supplemented by genuine aromatic-linker dispersion). One Mg-MOF-74 branch, 1c, also illustrates the silent-substitution defect class directly: its as-run framework Lennard-Jones parameters were plain UFF with a placeholder Mg \(\epsilon = 5.0\) K, not the ``reduced'' Becker values its label claimed --- a third audit-caught instance, retained here as a FAIL example rather than rebuilt. Polarisability, induction, and local coordination chemistry are plausible parts of the missing physics. A polarisable Mg-MOF-74 calculation is therefore a natural future target, but it is not a v0.4 result. Any such branch must be built as a new recipe, with its induced-dipole model, units, damping convention, and validation reference checked before it can affect a verdict \cite{Becker2018JPCC}.

The fourth missing piece is a quantitative consensus-reference layer. Some apparent failures depend as much on the reference observable as on the force field. A Henry constant from a low-pressure slope, a Toth fit, a Langmuir fit, a van't Hoff heat, and a calorimetric heat do not carry the same experimental meaning. Future versions should replace heuristic physical-accuracy bands with a systematic multi-source reference model where enough data exist. The goal is not to make the pass band wider, but to make the uncertainty in the reference observable explicit \cite{McCready2024JCTC}.

The fifth missing piece is a controlled route to machine-learned potentials. Machine-learned interatomic potentials may be useful where pairwise force fields miss the binding physics, especially at open metal sites. But they would change the recipe. A machine-learned potential is not a drop-in proof that Widom insertion is now predictive. It must be validated as its own energy model, with out-of-domain checks, hard-overlap handling, uncertainty diagnostics, and, where possible, comparison against electronic-structure reference calculations. No v0.4 strict pass used a machine-learned potential or GPU acceleration. Section~\ref{sec:v06} takes a first step on exactly this front: not a validated machine-learned adsorption branch, but a governance layer that characterises a machine-learned potential outside its training domain and, where the evidence supports it, screen-passes or withholds a verdict.

These missing pieces should not be confused with defects in the strict passes. The all-silica MFI + Ar (6c) and the genuine-UFF UiO-66 + CO\(_2\) (3a) results are valid for their locked recipes and references; the Si-CHA + CO\(_2\) branch (4c) is a cross-engine-checked Tier-B charged-zeolite control whose strict-versus-Tier-B verdict is engine-marginal. What is missing is the evidence needed to extend that reliability to harder classes of materials: flexible trapdoor zeolites, cationic zeolites, defective MOFs, and open-metal-site frameworks. The audit therefore points to a research programme rather than to a finished calculator.

These gaps fall into two groups that should not be confused. Section~\ref{sec:v06} reports one extension of exactly this governance logic that has already been carried out: cross-engine validation on CuspAI's kUPS engine and out-of-distribution governance of the machine-learned potential it ships. The remaining gaps named above --- flexible trapdoor methods, completed cationic-zeolite recipes, stronger open-metal-site models, and a quantitative consensus-reference layer --- are still future work, and each must be validated as a new locked recipe before it can affect any verdict.

The next version of the atlas should be judged by whether it closes these named gaps without weakening the standard of evidence. A new branch should not be promoted because it produces a convenient number. It should be promoted only when its structure, energy model, sampling method, reference observable, and verdict logic are explicit enough for another reader to reproduce and criticise. That is the standard v0.4 established, and it is the standard future work must preserve.
